\documentclass{standalone}
\usepackage{circuitikz}
\usetikzlibrary{calc}
\definecolor{circuitnotemark}{HTML}{FE00FE}
\begin{document}
\begin{circuitikz}[american, line width=0.8pt]
\ctikzset{bipoles/length=1.2cm}
\ctikzset{grounds/scale=1.36}
\foreach \x in {0,2,4,6,8,10,12} {\foreach \y in {0,-2} {\fill[gray, opacity=0.35] (\x,\y) circle (1.1pt);}}
\node[gray, font=\scriptsize] at (0,0.84) {1};
\node[gray, font=\scriptsize] at (2,0.84) {2};
\node[gray, font=\scriptsize] at (4,0.84) {3};
\node[gray, font=\scriptsize] at (6,0.84) {4};
\node[gray, font=\scriptsize] at (8,0.84) {5};
\node[gray, font=\scriptsize] at (10,0.84) {6};
\node[gray, font=\scriptsize] at (12,0.84) {7};
\node[gray, font=\scriptsize, anchor=east] at (-0.84,0) {a};
\node[gray, font=\scriptsize, anchor=east] at (-0.84,-2) {b};
\coordinate (b1) at (0,-2);
\coordinate (b3) at (4,-2);
\coordinate (b5) at (8,-2);
\coordinate (b7) at (12,-2);
\draw (b1) node[ocirc]{} node[above left]{IN}; % line 3
\node[ground] at (b3) {}; % line 4
\node[vcc] at (b5) {VCC}; % line 5
\node[vee] at (b7) {VEE}; % line 6
\path (-0.825,-1.393) rectangle (0.825,-0.207); % line 8
\node[anchor=west, inner sep=0, circuitnotemark, font=\large] at (0,-0.8) {X}; % line 8
\path (-1.397,-3.993) rectangle (1.397,-2.807); % line 9
\node[anchor=west, inner sep=0, circuitnotemark, font=\large] at (0,-3.4) {X}; % line 9
\path (2.508,-1.393) rectangle (5.492,-0.207); % line 10
\node[anchor=west, inner sep=0, circuitnotemark, font=\large] at (4,-0.8) {X}; % line 10
\path (2.349,-3.993) rectangle (5.651,-2.807); % line 11
\node[anchor=west, inner sep=0, circuitnotemark, font=\large] at (4,-3.4) {X}; % line 11
\path (6,-1.393) rectangle (10,-0.207); % line 12
\node[anchor=west, inner sep=0, circuitnotemark, font=\large] at (8,-0.8) {X}; % line 12
\path (6.603,-3.993) rectangle (9.397,-2.807); % line 13
\node[anchor=west, inner sep=0, circuitnotemark, font=\large] at (8,-3.4) {X}; % line 13
\path (10,-1.393) rectangle (14,-0.207); % line 14
\node[anchor=west, inner sep=0, circuitnotemark, font=\large] at (12,-0.8) {X}; % line 14
\path (10.603,-3.993) rectangle (13.397,-2.807); % line 15
\node[anchor=west, inner sep=0, circuitnotemark, font=\large] at (12,-3.4) {X}; % line 15
\node[anchor=south west, inner sep=2pt, circuitnotemark, font=\LARGE\bfseries] (circuittitle) at (current bounding box.north west) {X};
\path (circuittitle.south west) rectangle ++(6.736,0.853);
\end{circuitikz}
\end{document}
